\chapter{Categorical Catalysis: Aperture Mechanisms and Topological Networks}
\label{ch:catalysis}

\papernote{Source paper: \emph{Categorical Catalysis Mechanism} --- \texttt{docs/catalysis/}}

\begin{chapterabstract}
Catalysis is reinterpreted as categorical aperture reduction: an enzyme does not lower an activation energy barrier but reduces the categorical distance between reactant and product states by providing a structured aperture through which the completion trajectory passes. This interpretation unifies enzymatic catalysis, autocatalytic cycles, and heterogeneous catalysis (the Haber process) within the same categorical framework. The penultimate-step equilibrium theorem establishes why catalytic cycles are irreversible even when individual steps are reversible. Case studies: carbonic anhydrase II (the fastest known enzyme) and RuBisCO (multi-aperture catalysis in photosynthesis).
\end{chapterabstract}

\input{../docs/catalysis/sections/aperture-model.tex}
\input{../docs/catalysis/sections/autocatalytic-apertures.tex}
\input{../docs/catalysis/sections/gas-network-topology.tex}
\input{../docs/catalysis/sections/topological-network.tex}
\input{../docs/catalysis/sections/categorical-exclusion.tex}
\input{../docs/catalysis/sections/temporal-independence.tex}
\input{../docs/catalysis/sections/penultimate-step-equilibrium.tex}
\input{../docs/catalysis/sections/carbonic-anhydrase.tex}
\input{../docs/catalysis/sections/rubisco-multi-aperture.tex}
\input{../docs/catalysis/sections/haber-process.tex}
